\subsection{Verification Evidence and Rest-State Preservation}
\label{sec:numerical-verification}

The worked example uses the evidence standards from
Section~\ref{sec:numerical-methods} in three ways: manufactured-solution
verification checks the forced finite-volume operator, the geometry-rest test
checks preservation of the declared zero-flow equilibrium by the selected
comparison method, and admissibility diagnostics report the boundary, CFL,
positivity, and conservation regime of the accepted comparison runs. Secondary
checks for nonprincipal numerical surfaces remain in
Appendix~\ref{app:numerical-methods-details}.

\subsubsection{Evidence Hierarchy and Scope}
\label{subsec:verification-scope-convergence-stability}

The hierarchy is explicit. The strongest positive result is the
manufactured-solution test, because it compares the implemented forced operator
with a specified smooth exact solution and an independently assembled forcing
consistency check. The geometry-rest test then checks the named stationary
family used to initialize the comparison runs. The selected comparison method is
the geometry-rest-balanced finite-volume variant; the older MUSCL/Rusanov row
is retained as nonselected context because it does not exactly preserve the same
sampled rest state. Boundary, CFL, positivity, and conservation diagnostics then
state whether accepted runs remained numerically admissible. Self-convergence,
stationary-Stokes projection, broad rheology/profile sweeps, resolved-velocity
benchmark stages, and nonselected DG rows are secondary checks kept in
Appendix~\ref{app:numerical-methods-details}.

The section map has its own check because the 1D--3D comparison is only
meaningful after the resolved field is mapped to section area, physical flow,
and mean axial velocity. The synthetic check in
Table~\ref{tab:t1-cross-section-operator-validation} reports the workflow rows
for constant and affine node-centered axial fields. The package tests also check
closed-form single-tetrahedron and disjoint two-tetrahedron cuts with
independently computed transverse areas and affine means. These tests support
the plane-cut area, flow, and mean-velocity operator on controlled synthetic
geometry, not matched-data validation of the available 3D cases.

\begin{table}[!htb]
    \centering
    \scriptsize
    \caption{Synthetic verification check for the specified cross-section
    quadrature operator. Differences are measured against analytic constant or
    affine references on the synthetic tetrahedron; the
    plane-shift rows report changes relative to the $z=0.5$ affine cut.}
    \label{tab:t1-cross-section-operator-validation}
    \resizebox{\textwidth}{!}{%
    \IfFileExists{report/assets/tables/stenosis-comparison/cross_section_operator_validation.tex}{%
        % Generated by run_operator_validation; input mesh and fields are synthetic.
\begin{tabular}{lrrrrrrrl}
\toprule
case & z & shift & area & flow err. & mean err. & shift $\Delta A$ & shift $\Delta \bar u$ & status \\
\midrule
constant & 0.25 & 0.000e+00 & 0.2812 & 0.000e+00 & 0.000e+00 & -- & -- & pass \\
constant & 0.5 & 0.000e+00 & 0.125 & 0.000e+00 & 0.000e+00 & -- & -- & pass \\
constant & 0.75 & 0.000e+00 & 0.03125 & 0.000e+00 & 0.000e+00 & -- & -- & pass \\
affine & 0.25 & 0.000e+00 & 0.2812 & 0.000e+00 & 0.000e+00 & -- & -- & pass \\
affine & 0.5 & 0.000e+00 & 0.125 & 0.000e+00 & 0.000e+00 & -- & -- & pass \\
affine & 0.75 & 0.000e+00 & 0.03125 & 2.776e-17 & 8.882e-16 & -- & -- & pass \\
plane\_shift & 0.45 & -0.05 & 0.1513 & 0.000e+00 & 0.000e+00 & 0.02625 & -0.2167 & pass \\
plane\_shift & 0.5 & 0.000e+00 & 0.125 & 0.000e+00 & 0.000e+00 & 0.000e+00 & 0.000e+00 & pass \\
plane\_shift & 0.55 & 0.05 & 0.1012 & 1.110e-16 & 8.882e-16 & -0.02375 & 0.2167 & pass \\
\bottomrule
\end{tabular}
%
    }{%
        \begin{tabular}{ll}
        \toprule
        Status & Message \\
        \midrule
        missing & Operator-validation table unavailable in this source snapshot. \\
        \bottomrule
        \end{tabular}%
    }%
    }
\end{table}

Standard finite-volume context is supplied by conservative hyperbolic
discretization with consistent numerical fluxes and sources, CFL-controlled time
stepping, limiter and Rusanov dissipation choices, and SSP Runge--Kutta time
integration \cite{LeVeque2002FiniteVolumeMethodsForHyperbolicProblems}. In 1D
blood-flow practice, MUSCL/Rusanov and Lax--Friedrichs-type discretizations are
commonly assessed through exact-solution, benchmark, grid-refinement, and
equilibrium tests
\cite{Wang2014BloodFlowNetworks,BeckersKolbe2025LaxFriedrichsHemodynamics,BoileauEtAl2015Benchmark,BenemeritoEtAl2024OpenBF,LuccaEtAl2025FFRPrediction}.
Here those checks are empirical implementation evidence, not a formal
convergence theorem for the full nonlinear stenosis solver.

\subsubsection{Manufactured-Solution Evidence}
\label{subsec:mms-verification}

For verification only, the selected finite-volume operator is paired with a
manufactured forcing term. Production runs use no manufactured forcing. The
verification run supplies a smooth exact state
$(a^\star(z,t),q^\star(z,t))$, exact initial data, exact boundary traces, and
forcing residuals for the selected operator. In solver coordinates, let
\begin{flalign*}
    \quad
    s(z)&=\sin\!\left(\frac{\pi z}{L}\right),&
    a_0(z)&=R_0(z)^2,&
    \omega&=\frac{2\pi}{T_{\mathrm{mms}}}. &&
\end{flalign*}
The default manufactured state uses
\begin{flalign*}
    \quad
    \epsilon_a&=0.02,&
    U_{\mathrm{mms}}&=2.0\,\mathrm{cm}/\mathrm{s},&
    T_{\mathrm{mms}}&=0.02\,\mathrm{s}, &&
\end{flalign*}
and defines
\begin{flalign*}
    \quad
    a^\star(z,t)
    &=
    a_0(z)\left[1+\epsilon_a s(z)\cos(\omega t)\right], &&\\
    \quad
    q^\star(z,t)
    &=
    a_0(z)U_{\mathrm{mms}}s(z)^2\sin(\omega t). &&
\end{flalign*}
The tabulated run uses $0\le z\le L=6\,\mathrm{cm}$ and
$0\le t\le2\times10^{-3}\,\mathrm{s}$, equal to 10\% of the manufactured
period. The forcing terms are the residuals required by the selected flux and
source realization,
\begin{flalign*}
    \quad
    f_a(z,t)
    &=
    \pd_t a^\star+\pd_z F_a(a^\star,q^\star;z), &&\\
    \quad
    f_q(z,t)
    &=
    \pd_t q^\star+
    \pd_z F_q(a^\star,q^\star;z)-S_q(a^\star,q^\star;z), &&
\end{flalign*}
where $F_a,F_q$ and $S_q$ are the selected point flux and source functions.
The forcing consistency check uses separately expanded formulas rather than a
stored table replay, but it shares the manufactured states, geometry and
material parameters, and lower-level flux/source utilities used by the selected
operator. The table reports its maximum absolute mismatch. The run therefore
checks the implemented forced operator against a specified manufactured state; it
is implementation-verification evidence, not independent solver validation and
not a physical stenosis experiment.

The verification table reports cell-center discrete errors
\begin{flalign*}
    \quad
    \|e\|_1&=\frac{1}{N}\sum_i |e_i|,&
    \|e\|_2&=\left(\frac{1}{N}\sum_i e_i^2\right)^{1/2},&
    \|e\|_\infty&=\max_i |e_i|, &&
\end{flalign*}
with observed orders computed separately for each metric
$m\in\{1,2,\infty\}$ as
\begin{flalign*}
    \quad
    p_{\mathrm{obs}}^{(m)}
    &=
    \frac{\log(E_{\mathrm{coarse}}^{(m)}/E_{\mathrm{fine}}^{(m)})}
         {\log(\Delta_{\mathrm{coarse}}/\Delta_{\mathrm{fine}})}. &&
\end{flalign*}
These are discrete norm values for the stated cell-center error samples. The
$L^1$ metric tracks integrated error, the $L^2$ metric gives the RMS-like error
used in many convergence summaries, and the $L^\infty$ metric exposes the
largest pointwise discrepancy, including boundary, limiter, or source-term
localization. For spatial rows $\Delta$ is the grid spacing. For
requested-timestep rows the native solver also reports realized CFL diagnostics
because the accepted timestep is
$\min(\Delta t_{\mathrm{request}},\Delta t_{\mathrm{CFL}})$. The
timestep-insensitivity rows in the generated MMS table are therefore read
through accepted timestep and realized CFL, not as a clean temporal-order
study.

\IfFileExists{report/assets/tables/verification/mms_verification.tex}{%
    \begin{table}[!htb]
    \centering
    \scriptsize
    \caption{Manufactured-solution spatial verification errors in the cell-center discrete $L_1$, $L_2$, and $L_\infty$ metrics. The final column checks the inserted forcing against a separately expanded residual expression that still shares the manufactured states, parameters, and lower-level utilities.}
    \resizebox{\textwidth}{!}{%
    \begin{tabular}{@{}rrrrrrrrrr@{}}
        \toprule
        $N$ & $\Delta t_{\min}$ & CFL$_{\max}$ & $\|e_a\|_1$ & $\|e_a\|_2$ & $\|e_a\|_\infty$ & $\|e_q\|_1$ & $\|e_q\|_2$ & $\|e_q\|_\infty$ & forcing check \\
        \midrule
50 & 2.500e-06 & 0.02154 & 7.412e-07 & 1.018e-06 & 2.377e-06 & 0.001163 & 0.003447 & 0.017 & 0.000e+00 \\
100 & 2.500e-06 & 0.04309 & 1.895e-07 & 3.022e-07 & 9.223e-07 & 3.318e-04 & 0.001221 & 0.008493 & 0.000e+00 \\
200 & 2.500e-06 & 0.08617 & 4.866e-08 & 9.177e-08 & 3.661e-07 & 8.996e-05 & 4.310e-04 & 0.004241 & 0.000e+00 \\
400 & 2.500e-06 & 0.1723 & 1.245e-08 & 2.833e-08 & 1.457e-07 & 2.369e-05 & 1.520e-04 & 0.002119 & 0.000e+00 \\
800 & 2.500e-06 & 0.3447 & 3.183e-09 & 8.853e-09 & 5.923e-08 & 6.123e-06 & 5.362e-05 & 0.001059 & 0.000e+00 \\
1600 & 5.368e-07 & 0.45 & 8.104e-10 & 2.790e-09 & 2.410e-08 & 1.566e-06 & 1.892e-05 & 5.294e-04 & 0.000e+00 \\
3200 & 5.370e-07 & 0.45 & 2.058e-10 & 8.844e-10 & 9.881e-09 & 3.978e-07 & 6.675e-06 & 2.646e-04 & 0.000e+00 \\
        \bottomrule
    \end{tabular}%
    }
\end{table}

\begin{table}[!htb]
    \centering
    \scriptsize
    \caption{Observed spatial orders computed separately from adjacent-grid manufactured-solution errors in each discrete metric.}
    \begin{tabular}{@{}rrrrrrr@{}}
        \toprule
        grid pair & $p_{a,1}$ & $p_{a,2}$ & $p_{a,\infty}$ & $p_{q,1}$ & $p_{q,2}$ & $p_{q,\infty}$ \\
        \midrule
50--100 & 1.967 & 1.752 & 1.366 & 1.81 & 1.497 & 1.001 \\
100--200 & 1.962 & 1.719 & 1.333 & 1.883 & 1.502 & 1.002 \\
200--400 & 1.966 & 1.696 & 1.329 & 1.925 & 1.503 & 1.001 \\
400--800 & 1.968 & 1.678 & 1.299 & 1.952 & 1.503 & 1.001 \\
800--1600 & 1.974 & 1.666 & 1.297 & 1.967 & 1.503 & 1.0 \\
1600--3200 & 1.977 & 1.657 & 1.286 & 1.977 & 1.503 & 1.0 \\
        \bottomrule
    \end{tabular}
\end{table}

\begin{table}[!htb]
    \centering
    \scriptsize
    \caption{Manufactured-solution timestep-insensitivity rows on the finest MMS grid. These rows report accepted timestep and realized CFL rather than temporal order.}
    \resizebox{\textwidth}{!}{%
    \begin{tabular}{@{}rrrrrrrr@{}}
        \toprule
        $N$ & requested $\Delta t$ & accepted $\Delta t_{\min}$ & accepted $\Delta t_{\max}$ & CFL$_{\max}$ & $\|e_a\|_2$ & $\|e_q\|_2$ & forcing check \\
        \midrule
3200 & 2.000e-05 & 5.370e-07 & 8.165e-07 & 0.45 & 8.844e-10 & 6.675e-06 & 0.000e+00 \\
3200 & 1.000e-05 & 5.370e-07 & 8.165e-07 & 0.45 & 8.844e-10 & 6.675e-06 & 0.000e+00 \\
3200 & 5.000e-06 & 5.370e-07 & 8.165e-07 & 0.45 & 8.844e-10 & 6.675e-06 & 0.000e+00 \\
3200 & 2.500e-06 & 5.370e-07 & 8.165e-07 & 0.45 & 8.844e-10 & 6.675e-06 & 0.000e+00 \\
        \bottomrule
    \end{tabular}%
    }
\end{table}

}{%
    \begin{center}
    \small
    \emph{Manufactured-solution table unavailable in this source snapshot; see
    the appendix evidence summary.}
    \end{center}
}
The spatial rows provide bounded, positive implementation-verification evidence
for the specified forced operator through the finest regenerated MMS grid,
$N=3200$. Over this grid range, the area orders are about $1.96$--$1.98$ in
$L^1$, $1.66$--$1.75$ in $L^2$, and $1.29$--$1.37$ in $L^\infty$. The flow
orders are about $1.81$--$1.98$ in $L^1$, about $1.50$ in $L^2$, and about
$1.00$ in $L^\infty$. The maximum-norm columns therefore emphasize a more
localized part of the error than the averaged metrics. The evidence is read as
metric-specific implementation-verification behavior, not as a blanket
formal-order claim for every model variant.

\subsubsection{Geometry-Rest Preservation Test}
\label{subsec:rest-state-drift}

The continuous selected source-balanced law admits the geometry-rest state
$a_i=R_0(z_i)^2$, $q_i=0$ under compatible zero-flow boundary data
(Proposition~\ref{prop:selected-rest-state-compatibility}). A plain
MUSCL/Rusanov realization is not exactly well balanced for that state: its
area-equation viscosity sees the spatially varying $R_0$ profile and produces a
nonzero discrete residual unless the equilibrium is built into the discrete
operator. The selected comparison method therefore uses the finite-volume
geometry-rest-balanced variant specified in
Appendix~\ref{app:num-current-solver-operator}. It retains the MUSCL limiter
and Rusanov-type dissipation but applies the dissipative area jump to the
perturbation $a-a_0$, with $a_0(z_i)=R_0(z_i)^2$, and uses the matching
discrete wall source that cancels the elastic flux difference at the sampled
rest state. In zero-forcing, zero-inlet runs this preserves the declared
geometry-rest state to roundoff in focused package tests and in the same
rest-drift CSV/\TeX{} workflow used below.

Three levels of this statement must be kept separate. Continuous rest-state
compatibility is the cancellation property of the source-balanced differential
model. The instantaneous discrete residual at $t=0$ is the semidiscrete
right-hand side obtained after sampling that state and applying the selected
reconstruction, Rusanov flux, source, and boundary rules. The time-evolved rest
drift below is the accumulated response of the time integrator to that
residual. Table~\ref{tab:rest-state-residual-components} reports the initial
component magnitudes before the drift history is interpreted.

\IfFileExists{report/assets/tables/verification/rest_state_residual_components.tex}{%
    \begin{table}[!htb]
    \centering
    \scriptsize
    \caption{Initial geometry-rest residual component magnitudes at $t=0$. The area row is the Rusanov mass-flux residual. The momentum rows split the elastic-flux difference and wall/geometry source before summing them into the total flow residual. Values are maxima over cells in solver coordinates.}
    \label{tab:rest-state-residual-components}
    \resizebox{\textwidth}{!}{%
    \begin{tabular}{@{}lrrrrrrr@{}}
        \toprule
        Case & $N$ & $\max |R_a^{\mathrm{Rus}}|$ & $z_a$ & $\max |R_q^{\mathrm{el}}|$ & $\max |S_q^{\mathrm{wall}}|$ & $\max |R_q^{\mathrm{tot}}|$ & $z_{q,\mathrm{tot}}$ \\
        \midrule
C23 (22.56\%) & 400 & 0.000e+00 & 0.0075 & 6.799e+04 & 6.799e+04 & 0.000e+00 & 0.0075 \\
C23 (22.56\%) & 800 & 0.000e+00 & 0.00375 & 6.743e+04 & 6.743e+04 & 0.000e+00 & 0.00375 \\
C23 (22.56\%) & 1600 & 0.000e+00 & 0.001875 & 6.737e+04 & 6.737e+04 & 0.000e+00 & 0.001875 \\
C23 (22.56\%) & 3200 & 0.000e+00 & 9.375e-04 & 6.738e+04 & 6.738e+04 & 0.000e+00 & 9.375e-04 \\
C40 & 400 & 0.000e+00 & 0.0075 & 1.043e+05 & 1.043e+05 & 0.000e+00 & 0.0075 \\
C40 & 800 & 0.000e+00 & 0.00375 & 1.043e+05 & 1.043e+05 & 0.000e+00 & 0.00375 \\
C40 & 1600 & 0.000e+00 & 0.001875 & 1.042e+05 & 1.042e+05 & 2.068e-25 & 1.761 \\
C40 & 3200 & 0.000e+00 & 9.375e-04 & 1.042e+05 & 1.042e+05 & 4.136e-25 & 1.763 \\
        \bottomrule
    \end{tabular}%
    }
\end{table}

}{%
    \begin{center}
    \small
    \emph{Rest-state residual-component table unavailable in this source
    snapshot; see the appendix evidence summary.}
    \end{center}
}

The geometry-rest test is a numerical-methods result in its own right. The
continuous source-balanced model admits the quiescent stenotic geometry as an
exact rest state, and the selected comparison realization is required to
preserve that sampled state on the production grid. The zero-forcing rest test
starts the selected C23 (22.56\%) and 40\% stenosis geometries from the rest
state with a zero-velocity inlet and reports
\begin{flalign*}
    \quad
    &\max_i |q_i|,&
    &\max_i |a_i-R_0(z_i)^2|,&
    &\left|\sum_i a_i\Delta z-\sum_i a_i(0)\Delta z\right| &&
\end{flalign*}
as functions of grid size and elapsed time. This test fixes the interpretation
of any comparison run initialized from the same geometry-rest state. Its
terminal elapsed time is the diagnostic endpoint $t=1.0\,\mathrm{s}$; this is
not the imported resolved-data matching time $0.9995\,\mathrm{s}$ used later for
the final-time velocity comparison.

\IfFileExists{report/assets/tables/verification/rest_state_drift.tex}{%
    \begin{center}
    \fbox{\begin{minipage}{0.92\linewidth}
    \small
    \textbf{Rest-preservation check.} The zero-forcing geometry-rest runs use
    the same geometry-rest-balanced finite-volume method as the refreshed
    comparison record. The generated rows report near-roundoff solver flow,
    area drift, and finite-volume balance residuals for the declared rest
    family.
    \end{minipage}}
    \end{center}
    \begin{table}[!htb]
    \centering
    \scriptsize
    \caption{Zero-forcing, zero-inlet geometry-rest drift summary. The table reports the combined requested/applied inlet-flow status, the largest cell flow over positive elapsed times, and the final balance residual; the full diagnostic table remains in Appendix~\ref{app:num-full-rest-state-drift-grid}.}
    \resizebox{\textwidth}{!}{%
    \begin{tabular}{@{}rrlrrrrr@{}}
        \toprule
        Severity & $N$ & inlet $q_{\mathrm{in}}$ req./appl. & $t_{\mathrm{peak}}$ & peak $\max |q_i|$ & $z_{\max |q|}$ & final $\max |q_i|$ & final balance residual \\
        \midrule
23 & 400 & $(\mathrm{0.000e+00})/(\mathrm{0.000e+00})$ & 0.001 & 0.04387 & 2.078 & 0.03839 & 0.000e+00 \\
23 & 800 & $(\mathrm{0.000e+00})/(\mathrm{0.000e+00})$ & 0.001 & 0.01116 & 2.074 & 0.00973 & 0.000e+00 \\
23 & 1600 & $(\mathrm{0.000e+00})/(\mathrm{0.000e+00})$ & 0.001 & 0.002807 & 2.072 & 0.002442 & 0.000e+00 \\
23 & 3200 & $(\mathrm{0.000e+00})/(\mathrm{0.000e+00})$ & 0.001 & 7.030e-04 & 2.071 & 6.110e-04 & 0.000e+00 \\
40 & 400 & $(\mathrm{0.000e+00})/(\mathrm{0.000e+00})$ & 0.001 & 0.07053 & 2.078 & 0.06339 & 0.000e+00 \\
40 & 800 & $(\mathrm{0.000e+00})/(\mathrm{0.000e+00})$ & 0.001 & 0.01782 & 2.074 & 0.01605 & 0.000e+00 \\
40 & 1600 & $(\mathrm{0.000e+00})/(\mathrm{0.000e+00})$ & 0.001 & 0.004465 & 2.072 & 0.004028 & 0.000e+00 \\
40 & 3200 & $(\mathrm{0.000e+00})/(\mathrm{0.000e+00})$ & 0.001 & 0.001116 & 2.071 & 0.001009 & 0.000e+00 \\
        \bottomrule
    \end{tabular}%
    }
\end{table}

}{%
    \begin{center}
    \small
    \emph{Rest-state drift table unavailable in this source snapshot; see
    the appendix evidence summary.}
    \end{center}
}

The accepted zero-forcing rest-state run has zero requested and applied inlet
flow, and the finite-volume balance
$\Delta\!\int a\,dz+\int(\widehat F^a_{\mathrm{out}}-\widehat F^a_{\mathrm{in}})\,dt$
closes to near roundoff in the produced CSV. The production comparison-flow
scale is
$q_{\mathrm{comp}}\approx2.292/\pi=0.7296\,\mathrm{cm}^3/\mathrm{s}$ in the
solver coordinate, corresponding to
$Q_{\mathrm{comp}}\approx2.292\,\mathrm{cm}^3/\mathrm{s}$ in physical circular
flow, as used in the axial-flow comparison
(Figure~\ref{fig:t1-axial-flow-comparison}). On the $N=400$ production
comparison grid, the final zero-inlet rest-flow values are roundoff-scale
relative to this comparison-flow scale. The older non-well-balanced
MUSCL/Rusanov drift row is therefore retained as a historical method diagnostic,
not as the governing limitation for the reported comparison runs.

An additional method-sensitivity check asks whether rest preservation is tied to
the selected geometry-rest-balanced construction rather than to every available
finite-volume surface. Table~\ref{tab:rest-method-sensitivity} runs the same
zero-inlet geometry-rest diagnostic at $N=400$ through historical
finite-volume methods, the selected \texttt{fv-wb-geometry-rest} method, and the
50\% stenosis geometry as a higher-severity stress case. First-order fluxes are
more diffusive, and WENO3 reduces the instantaneous residual and drift by more
than an order of magnitude, but those historical executable rows still do not
preserve the sampled geometry-rest state exactly. The Lax--Wendroff row is also
diagnostic context: it reduces the drift relative to the plain MUSCL row for
this scratch grid, but it is not the selected comparison method and does not
preserve the declared rest state to roundoff. The
\texttt{fv-wb-geometry-rest} rows are the rows used for the refreshed
comparison record.

\begin{table}[!htb]
    \centering
    \scriptsize
    \caption{Method-sensitivity check for the zero-inlet geometry-rest
    diagnostic at $N=400$. The final $\max |q|$ column reports the largest
    sampled solver-flow magnitude at the final requested diagnostic time,
    $t=0.1$ s, for executable rows. Values are solver-flow units cm$^3$/s
    except for the
    residual column, which reports the initial discrete flow residual
    magnitude.}
    \label{tab:rest-method-sensitivity}
    \resizebox{\textwidth}{!}{%
    \begin{tabular}{@{}llrrrrl@{}}
\toprule
Method & Case & $\max |R_q(0)|$ & $t_{\mathrm{peak}}$ & peak $\max |q|$ & final $\max |q|$ & Status \\
\midrule
FV first-order & C23 & 526.2 & 0.001 & 0.6043 & 0.5367 & ok \\
FV first-order & C40 & 917 & 0.001 & 1.012 & 0.86 & ok \\
FV first-order & C50 & 1.167e+03 & 0.001 & 1.224 & 1.028 & ok \\
FV MUSCL & C23 & 630.3 & 0.001 & 0.04387 & 0.03992 & ok \\
FV MUSCL & C40 & 1.437e+03 & 0.001 & 0.07053 & 0.06172 & ok \\
FV MUSCL & C50 & 1.397e+03 & 0.001 & 0.08193 & 0.07736 & ok \\
FV WENO3 & C23 & 31.42 & 0.1 & 0.003972 & 0.003972 & ok \\
FV WENO3 & C40 & 89.18 & 0.001 & 0.008444 & 0.008242 & ok \\
FV WENO3 & C50 & 136.7 & 0.001 & 0.01186 & 0.01185 & ok \\
FV Lax--Wendroff & C23 & -- & -- & -- & -- & not available \\
FV Lax--Wendroff & C40 & -- & -- & -- & -- & not available \\
FV Lax--Wendroff & C50 & -- & -- & -- & -- & not available \\
\bottomrule
\end{tabular}

    }
\end{table}

\begin{figure}[!htb]
    \centering
    \figtikz{rest-state-flow-profiles}
    \caption{Zero-inlet rest-state solver-flow profiles on the $N=800$ grid for
    the selected geometry-rest-balanced finite-volume method. The profiles
    remain at roundoff scale under the declared rest-family diagnostic.}
    \label{fig:rest-state-flow-profiles}
\end{figure}

This geometry-rest result does not turn the later section-mean velocity
comparison into external validation, but it removes the previous production-grid
rest-drift confound from that comparison record. The reported 1D--3D
discrepancy remains a conditional diagnostic comparison, conditioned on the
specified model, boundary approximation, observation rule, and available
resolved data, rather than a clean prediction-error estimate.

\subsubsection{Boundary, CFL, Positivity, and Conservation Diagnostics}
\label{subsec:boundary-diagnostics}

The accepted runs report timestep, realized CFL, characteristic-speed,
solver-volume, positivity, and finite-volume balance diagnostics. The
fixed-area characteristic boundary rule is used under the subcritical sign
condition $\lambda_-<0<\lambda_+$. A positive characteristic radicand alone is
not the boundary-regime evidence; the persisted sign extrema and
subcriticality margin are the relevant diagnostics.

In the rest-state table, the finest-grid C23 (22.56\%) and 40\% stenosis
$t=1$ s rows retain positive subcritical margins, zero requested and applied
inlet flow, signed solver-volume defects, boundary-flux integrals, and
near-roundoff finite-volume balance residuals. The rest CSV reports zero
positivity projection counts for these rest tests. The production comparison
diagnostics in
Table~\ref{tab:t1-alpha-radicand-diagnostics} provide the corresponding
characteristic-speed evidence for the runs later used in the velocity-output
comparison. These diagnostics are admissibility and finite-volume balance
checks. They do not replace the separate geometry-rest preservation test above.

\begin{table}[!htb]
    \centering
    \scriptsize
    \caption{Subcritical-boundary diagnostics for the same two 1D runs used in
    the resolved-velocity comparison. The characteristic-radicand values are in
    solver units of cm$^2$/s$^2$; the signs of $\lambda_-$ and $\lambda_+$
    support the boundary-regime gate in
    Assumption~\ref{ass:1d-subcritical-boundary-regime}.}
    \label{tab:t1-alpha-radicand-diagnostics}
    \resizebox{\textwidth}{!}{%
    \begin{tabular}{@{}lrrrrrrrr@{}}
        \toprule
        \textbf{Case}
            & $\boldsymbol{\min \alpha_{\mathrm{eff}}}$
            & $\boldsymbol{\max \alpha_{\mathrm{eff}}}$
            & $\boldsymbol{\min \chi^2}$
            & $\boldsymbol{\min\lambda_-}$
            & $\boldsymbol{\max\lambda_-}$
            & $\boldsymbol{\min\lambda_+}$
            & $\boldsymbol{\max\lambda_+}$
            & $\boldsymbol{\min\{\lambda_+,-\lambda_-\}}$ \\
        \midrule
        C23 (22.56\%) stenosis & 1.33068 & 1.33333 & $8.20\times 10^5$ & -1017.08 & -855.19 & 955.53 & 1058.85 & 855.19 \\
        40\% stenosis & 1.32500 & 1.33333 & $6.36\times 10^5$ & -1017.22 & -713.53 & 881.03 & 1058.97 & 713.53 \\
        \bottomrule
    \end{tabular}
    }
\end{table}

\begin{table}[!htb]
    \centering
    \scriptsize
    \pgfplotstableread[col sep=space]{report/assets/data/stenosis-comparison/production-diagnostics-reference.dat}\TOneProductionDiagnostics
    \caption{Production diagnostics from the C23 (22.56\%) and 40\% severity 1D
    comparison runs completed at $T=0.9995$ s. These rows summarize timestep,
    CFL, positivity, area, and finite-volume balance for the 1D runs used in
    the resolved-velocity comparison.}
    \label{tab:t1-production-diagnostics}
    \resizebox{\textwidth}{!}{%
    \pgfplotstabletypeset[
        columns={
            case,
            spatial_method,
            target_time,
            completed_time,
            cross_model_time_offset,
            dt_min,
            dt_max,
            cfl_max,
            min_Aphys,
            volume_defect,
            positivity_count,
            area_flux_balance,
            rhs_area_max,
            rhs_flow_max
        },
        every head row/.style={before row=\toprule, after row=\midrule},
        every last row/.style={after row=\bottomrule},
        columns/case/.style={string type,column type=l,column name=\textbf{Stenosis}},
        columns/spatial_method/.style={string type,column type=l,column name=\textbf{Method}},
        columns/target_time/.style={fixed,precision=4,column name=$\boldsymbol{T_{\mathrm{target}}}$},
        columns/completed_time/.style={fixed,precision=4,column name=$\boldsymbol{T_{\mathrm{done}}}$},
        columns/cross_model_time_offset/.style={sci,precision=2,column name=$\boldsymbol{|\Delta t_{\mathrm{1D--3D}}|}$},
        columns/dt_min/.style={sci,precision=2,column name=$\boldsymbol{\Delta t_{\min}}$},
        columns/dt_max/.style={sci,precision=2,column name=$\boldsymbol{\Delta t_{\max}}$},
        columns/cfl_max/.style={fixed,precision=2,column name=\textbf{CFL$_{\max}$}},
        columns/min_Aphys/.style={fixed,precision=5,column name=$\boldsymbol{\min \pi a}$},
        columns/volume_defect/.style={sci,precision=2,column name=$\boldsymbol{\Delta\!\int a\,dz}$},
        columns/positivity_count/.style={int detect,column name=\textbf{pos. count}},
        columns/area_flux_balance/.style={sci,precision=2,column name=\textbf{area-flux balance}},
        columns/rhs_area_max/.style={sci,precision=2,column name=$\boldsymbol{\max|\dot a|}$},
        columns/rhs_flow_max/.style={fixed,precision=5,column name=$\boldsymbol{\max|\dot q|}$}
    ]\TOneProductionDiagnostics
    }
\end{table}

\subsubsection{Secondary Numerical Checks}
\label{subsec:verification-summary}

The fourth evidence tier is kept out of the main argument.
Appendix~\ref{app:numerical-methods-details} retains package-level checks,
sampled self-convergence, nonselected DG p/h behavior, resolved-velocity
benchmark context, rheology/profile sensitivity, and fixed-wall
stationary-Stokes refinement. Those secondary checks can reveal runnability,
boundedness, solver-policy, or sensitivity issues. They do not replace the MMS
evidence, and they do not change the geometry-rest preservation conclusion.

\FloatBarrier
